\section{What these runs can and cannot answer}
\label{sec:limits}

Everything in this section comes from the sixteen runs the analyzer used on the last
run of \texttt{limodel analyze}, and every number was read off the code rather than
recalled. Nothing here is a criticism of the code. The code is honest about what it
was handed; this section is about what was handed to it.

Twenty-one runs are in \texttt{data/runs.csv}. Sixteen reach the fitting. The five that
do not are dealt with in \S\ref{sec:limits:waiting}.

\subsection{The design is a cross, not a grid}
\label{sec:limits:cross}

A second-order polynomial in two settings wants a grid: every current density crossed
with every feed ratio. These runs are not that. They are two arms meeting at one corner.

The first arm moves current density with the feed ratio pinned at its lowest value,
$R = 1.43\unit{mol/mol}$. Table~\ref{tab:limits:current-arm} is every run on it.

\begin{table}[htbp]
  \centering
  \caption{The current-density arm. Feed ratio $1.43\unit{mol/mol}$ on all thirteen runs.
    Current density is the median supply current of the run over the membrane area, which
    is why the levels are not round numbers.}
  \label{tab:limits:current-arm}
  \begin{tabular}{l l l}
    \toprule
    $J$ ($\unit{mA/cm^2}$) & Membrane & Runs \\
    \midrule
    4.50  & GO-10            & AJ \\
          & Pristine         & AM \\
    7.50  & GO-10            & AH, AI \\
          & Pristine         & AL \\
          & PDA/PEI/15CE/BT  & AR \\
          & PDA-GO10/PEI/BT  & AS \\
    13.67 & GO-10            & AE \\
          & Pristine         & AK \\
    25.67 & GO-10            & AN, AP \\
          & Pristine         & AO, AQ \\
    \bottomrule
  \end{tabular}
\end{table}

The second arm moves the feed ratio with the current density pinned at
$13.67\unit{mA/cm^2}$, and it exists on one membrane only.

\begin{table}[htbp]
  \centering
  \caption{The feed-ratio arm. Current density $13.67\unit{mA/cm^2}$, membrane GO-10, on
    all four runs. AE is the corner the two arms share.}
  \label{tab:limits:ratio-arm}
  \begin{tabular}{l l l}
    \toprule
    $R$ ($\unit{mol/mol}$) & Membrane & Run \\
    \midrule
    1.43  & GO-10 & AE \\
    2.86  & GO-10 & AF \\
    5.71  & GO-10 & AG \\
    11.42 & GO-10 & AT \\
    \bottomrule
  \end{tabular}
\end{table}

Put the two tables side by side and the shape is plain. Four current-density levels exist
and four feed-ratio levels exist, so a full grid would hold sixteen cells. Seven are
occupied. Every run off the corner sits on one arm or the other, never in the middle of
the field.

This is why \texttt{SecondOrderPolynomial} has no $x_1x_2$ term. The code says so in a
comment beside \texttt{TERMS\_IN\_ORDER}, and \texttt{CANDIDATE\_TERMS} in
\texttt{fitted\_outcome.py} never offers one. An interaction term answers the question
\emph{does the best current density depend on how much magnesium is in the feed}. To
answer it, the feed ratio has to move at more than one current density. In these runs it
moves at one. The term is not missing from the code because it was forgotten. It is
missing because there is nothing in the data for it to be fitted to. That is a fact about
the experiments, and no change to the code can repair it.

There is a second thing the tables do not show, and it matters more than the missing
interaction. The four runs at $25.67\unit{mA/cm^2}$ --- AN, AO, AP and AQ --- all started
their concentrate at $5\unit{g/L}$. Every other run started it at $0.5\unit{g/L}$. AN and
AO started from sodium chloride; AP and AQ from lithium chloride. So the top level of the
current-density arm changes the concentrate at the same time as it changes the current,
and \texttt{Result} carries neither \texttt{conc\_start\_salt} nor
\texttt{conc\_start\_g\_l}, so the fitting cannot see the change and cannot allow for it.
Any curvature the fit finds at the top of the current range is a current effect and a
concentrate effect added together, with no way to tell the two apart.

\subsection{What the bench gets when it repeats itself}
\label{sec:limits:replicates}

\texttt{ReplicateScatter} groups runs by membrane and by both settings rounded to two
places, then pools the scatter inside the groups. Three groups have more than one run.
Table~\ref{tab:limits:replicates} gives them with the gap between the two members.

\begin{table}[htbp]
  \centering
  \caption{Every repeated design point among the sixteen runs, and the gap between the
    pair in each of the three fitted outcomes. Feed ratio is $1.43\unit{mol/mol}$ on all
    six runs.}
  \label{tab:limits:replicates}
  \begin{tabular}{l l r r r r}
    \toprule
    & & $J$ & Recovery & & Voltage \\
    Pair & Membrane & ($\unit{mA/cm^2}$) & (points) & $P$ & ($\unit{V}$) \\
    \midrule
    AH, AI & GO-10    & 7.50  & 1.07  & 0.356 & 0.33 \\
    AN, AP & GO-10    & 25.67 & 13.18 & 0.596 & 1.63 \\
    AO, AQ & Pristine & 25.67 & 3.55  & 0.087 & 0.20 \\
    \bottomrule
  \end{tabular}
\end{table}

Pooled, GO-10's two repeats give a standard deviation of 6.61 points of recovery, 0.347
in permselectivity and $0.832\unit{V}$, each on two degrees of freedom. Pristine's one
repeat gives 2.51 points, 0.062 and $0.141\unit{V}$ on one degree of freedom. These are
the numbers the report prints as the bench scatter.

Now put the largest gap next to the effect the fit is chasing. Along the current-density
arm, GO-10's recovery reads 21.1\% at $4.50\unit{mA/cm^2}$, 36.6\% and 37.7\% at 7.50,
38.3\% at 13.67, and 39.5\% and 26.3\% at 25.67. Averaging the pairs, the step from 7.50
to 13.67 is 1.2 points and the step from 13.67 to 25.67 is $-5.4$ points. The AN--AP gap
is 13.2 points. It is larger than both of those steps put together, and larger than the
whole span of the level averages above $4.50\unit{mA/cm^2}$. Only the first step, 4.50 to
7.50 at 16.0 points, is bigger than the scatter at a single point.

Permselectivity is no better placed. The AN--AP gap of 0.596 is twice the 0.288 that
separates the 7.50 and 13.67 levels, and most of the 0.878 that separates 13.67 from
25.67. Voltage is the one outcome that survives the comparison: the arm climbs
$6.59 \to 10.39 \to 16.0 \to 17.73\unit{V}$ and every step is larger than the
$1.63\unit{V}$ gap.

So for two of the three outcomes the bench's own scatter at one setting is as large as the
difference between settings. The report says the same thing in the language of error:
GO-10's leave-one-out error is 0.561 against a scatter of 0.347 in permselectivity and
7.12 points against 6.61 in recovery, and in both cases it prints \emph{this fit has
earned no conclusion}. No amount of fitting fixes this. Until a repeat at a given setting
lands closer than the difference between settings, a curve drawn through these points
cannot be told apart from a curve drawn through the scatter.

One caveat about that 13.2-point gap, because it is the number doing most of the damage.
AN and AP are a repeat only in the two columns the grouping looks at. They differ in the
concentrate they started from, sodium chloride against lithium chloride, both at
$5\unit{g/L}$. The same is true of AO and AQ. AH and AI are the only pair in the data
whose set-up columns all agree. So part of what is being charged to bench scatter may be
a real effect of the starting concentrate --- which would be better news for the
repeatability and worse news for the design, since it means the top current level is
confounded twice over.

\subsection{Two membranes cannot be fitted at all}
\label{sec:limits:membranes}

Of the sixteen runs, nine are GO-10 and five are Pristine. The other two are a single
run each of PDA/PEI/15CE/BT (AR) and of PDA-GO10/PEI/BT (AS).
\texttt{Surrogate.FEWEST\_RUNS} is 3: an intercept and one
slope are two coefficients, and leave-one-out needs a row to spare. One run is not three,
so the report prints

\begin{quote}
\texttt{not fitted}\\
\texttt{PDA/PEI/15CE/BT has too few runs to fit,}\\
\texttt{1 against the 3 a fit needs}
\end{quote}

and the same for PDA-GO10/PEI/BT. Both of those runs sit at the same point, GO-10's
best-repeated setting of $7.50\unit{mA/cm^2}$ and $R = 1.43$. They give a permselectivity
of 1.43 and 1.23 against GO-10's 2.39 and 2.03 and Pristine's 0.28 at that same point, so
they are a reading, but one reading each. Nothing can be said about how either coating
responds to current density or to feed ratio, because neither was ever moved on them.

Pristine is fitted, but only in one direction. Its five runs never leave
$R = 1.43\unit{mol/mol}$, so \texttt{DesignPoints.find\_dropped\_settings} reports the
feed ratio as dropped and the report says so in words: \emph{fitted in current density,
with the feed ratio held at one value across these runs and never tested}. A Pristine
prediction at a higher feed ratio is not a cautious extrapolation. It is a number the
runs contain no information about at all.

\subsection{One polynomial over all sixteen runs}
\label{sec:limits:pooled}

The analyzer fits each membrane separately. The obvious alternative is one polynomial
over everything, which is what \texttt{docs/polynomial.md} describes. It is worth running,
because what comes back is the clearest statement of the problem.

The fit uses the project's own machinery: \texttt{DesignPoints} to code the settings and
\texttt{SecondOrderPolynomial} with all five of \texttt{TERMS\_IN\_ORDER}, forced in
rather than grown term by term. The coded settings are logged, then centred and scaled so
the lowest run sits at $-1$ and the highest at $+1$:

\begin{align}
  x_1 &= \frac{\ln\!\left(J / 10.747\right)}{0.87056}, \label{eq:limits:x1} \\
  x_2 &= \frac{\ln\!\left(R / 4.0381\right)}{1.03972}. \label{eq:limits:x2}
\end{align}

So $x_1 = 0$ is $10.75\unit{mA/cm^2}$, one coded unit is a factor of 2.39 in current
density, and one coded unit in $x_2$ is a factor of 2.83 in feed ratio. Fitting lithium
recovery in percent over all sixteen runs gives

\begin{equation}
  \widehat{\text{recovery}} = 12.09 + 2.03\,x_1 - 12.80\,x_2 - 8.97\,x_1^2 + 2.55\,x_2^2 .
  \label{eq:limits:pooled}
\end{equation}

Three numbers judge it:

\begin{itemize}[nosep]
  \item $R^2 = 0.299$. The surface accounts for 30\% of the variation in recovery across
        the sixteen runs.
  \item Leave-one-out error, through \texttt{LeaveOneOutError}, is 13.98 points.
  \item The standard deviation of the sixteen recoveries themselves is 12.96 points.
\end{itemize}

The leave-one-out error is the larger. That has a plain meaning. Hold a run back, fit
\eqref{eq:limits:pooled} to the other fifteen, predict the held-back run, and the
prediction misses by more than you would have missed by ignoring the settings entirely
and quoting the average recovery. The comparison can be made exactly: fitting an
intercept alone under the same leave-one-out loop gives 13.39 points. So the four extra
terms make out-of-sample prediction worse, not better. This is why
\texttt{FittedOutcome.find\_best\_polynomial} grows a model one term at a time and stops
when the next term does not lower the leave-one-out error --- a five-term fit forced onto
sixteen runs spread over seven design cells is fitting the scatter.

The surface's own peak says the same thing from another direction. Searched over the
coded box the runs cover, \eqref{eq:limits:pooled} tops out at 27.6\% recovery, at
$J = 11.9\unit{mA/cm^2}$ and $R = 1.43\unit{mol/mol}$. Four runs already measured beat it:
AE at 38.3\%, AH at 36.6\%, AI at 37.7\% and AN at 39.5\%, all GO-10 at $R = 1.43$. A
best-of-surface that sits more than ten points below four measurements in hand is not an
optimum. It is an average of membranes that behave nothing alike, and it is also why the
per-membrane report ends with \emph{no best setup: nothing inside the measured range
cleared the lithium recovery floor}. GO-10's own fitted recovery peaks at 33.3\%, under
the $35\%$ floor in \texttt{BestSetupUnderTheGates.RECOVERY\_FLOOR\_PCT}, while the same
four runs measured above it.

\texttt{docs/polynomial.md} allows the membrane in as a 0/1 indicator rather than as a
separate fit per membrane. Adding one such column, 1 for GO-10 and 0 for everything else,
to the same five terms gives

\begin{equation}
  \widehat{\text{recovery}} = -8.59 + 1.61\,x_1 - 18.13\,x_2 - 10.16\,x_1^2
    + 9.21\,x_2^2 + 20.15\,g ,
  \label{eq:limits:indicator}
\end{equation}

where $g = 1$ on a GO-10 run. $R^2$ rises from 0.299 to 0.819 and the leave-one-out error
falls from 13.98 to 7.65 points, now comfortably under the 12.96-point standard deviation.
One binary column explains more than all four setting terms together, and its coefficient
is 20 points of recovery --- larger than anything either setting does across its whole
measured range.

Read that carefully. It does not mean \eqref{eq:limits:indicator} is a good model. It
means the membrane is the dominant variable in this data set and the two dialled settings
are secondary to it. The indicator fit also still leans on the Pristine runs to pin
$x_2^2$, and Pristine never left $R = 1.43$, so the $+9.21$ on $x_2^2$ is carried by GO-10
alone. The per-membrane approach the analyzer actually takes is the more honest reading of
the same fact: if the membrane is worth 20 points, fit the membranes apart.

\subsection{The runs waiting on the lab}
\label{sec:limits:waiting}

Five of the twenty-one runs are dropped, all for the same reason, which the journal
records as \emph{no lab results at both ends}. Without a sample at the first and last
minute there is no amount that crossed, so there is no recovery and no permselectivity to
fit. Table~\ref{tab:limits:waiting} is the five, with the settings read from their own
rows and power logs.

\begin{table}[htbp]
  \centering
  \caption{The dropped runs and the settings they were made at. Current density is the
    median of each power log over the $6\unit{cm^2}$ area; all five come out at
    $13.5\unit{mA/cm^2}$.}
  \label{tab:limits:waiting}
  \begin{tabular}{l l r r}
    \toprule
    Run & Membrane & $J$ ($\unit{mA/cm^2}$) & $R$ ($\unit{mol/mol}$) \\
    \midrule
    AX & Pristine & 13.5 & 2.86 \\
    AY & Pristine & 13.5 & 5.71 \\
    AU & Pristine & 13.5 & 11.42 \\
    AV & Pristine & 13.5 & 17.13 \\
    AW & GO10     & 13.5 & 17.13 \\
    \bottomrule
  \end{tabular}
\end{table}

Four of the five are a Pristine feed-ratio sweep at one current density: 2.86, 5.71, 11.42
and 17.13, the same ladder GO-10 already has, with one rung added past its top. That is
exactly the dimension \S\ref{sec:limits:cross} found missing --- Pristine currently has no
feed-ratio movement at all, which is why the report calls the feed ratio dropped and
never tested on it. The fifth, AW, extends GO-10's own sweep to 17.13.

Their current density, $13.5\unit{mA/cm^2}$, is the same nominal setting as the existing
feed-ratio arm at $13.67\unit{mA/cm^2}$; the $0.17$ difference is the median of one power
log against another, not a different dial position. The concentrate start on all five is
lithium chloride at $0.5\unit{g/L}$, matching the existing arm, so these runs do not carry
the confound described in \S\ref{sec:limits:cross}.

Two things about them are worth knowing before the results land. First, the sweep is still
a sweep: it adds a second arm at the same current density, so after it arrives the feed
ratio will move on two membranes but still at one current density, and there will still be
no support for an interaction term. Second, AW's membrane is spelled \texttt{GO10} in
\texttt{data/runs.csv} while the nine fitted runs are spelled \texttt{GO-10}. The field is
a plain string --- \texttt{cem} in \texttt{models/run.py} and in the schema --- and
\texttt{Surrogate.find\_membranes} groups on it exactly as written, so AW will arrive as a
fourth membrane with one run and be refused for having too few runs, unless the spelling
is made to match.

\subsection{The balance gaps, and what they mean for quoting any number}
\label{sec:limits:gaps}

\texttt{BalanceGap} reports, as a percentage, how much of what the dilute lost was never
found in the concentrate and the rinse. It is a diagnostic. Nothing is ever corrected by
it. But it is the best single check on whether the lab numbers a run was built from hang
together, so it belongs in any judgement of how far these results can be pushed.

Over the sixteen runs the lithium gaps run from $-272.5\%$ to $+71.3\%$, with a median of
$+8.5\%$. Ten are positive, six negative. Seven of the sixteen are more than 25\% away
from zero in either direction, and only five are inside $\pm 10\%$. The magnesium gaps are
worse: $-76.8\%$ to $+67.6\%$, median $+17.5\%$, ten of sixteen outside $\pm 25\%$ and
only three inside $\pm 10\%$.

Both signs have a meaning, and only one of them is benign.

A positive gap means the dilute lost more than the other two tanks gained. That is what
hold-up looks like: liquid left in sample bottles and lines. \texttt{BalanceGap} counts an
unweighed rinse as zero gain, which it explains as the honest way for a diagnostic to be
wrong, and that choice pushes gaps positive. So a modest positive gap is roughly what a
sound run should show, and the median of $+8.5\%$ in lithium is consistent with that.

A negative gap means the concentrate and rinse together gained more metal than the dilute
lost. As mass, that cannot happen. Six runs show it in lithium and six in magnesium, so in
those runs something in the concentrations, the dilution factors or the tank weights does
not reconcile, and the amounts that crossed --- the numerator of both permselectivity and
recovery --- inherit that.

Two further points about size. The gap is a ratio, so it blows up when the denominator is
small. AT's $-272.5\%$ is the clearest case: its dilute lost $5.1\unit{mg}$ of lithium
over the whole run while the concentrate gained $18.8\unit{mg}$. A $14\unit{mg}$
discrepancy on a $5\unit{mg}$ loss is a huge percentage and a small absolute error, and AT
is the highest feed ratio in the data, the far end of the feed-ratio arm. The point at the
end of the lever is the one with the least lithium movement behind it. The other is that
the magnesium gaps are consistently larger than the lithium ones, which bears directly on
permselectivity, since $P$ divides a lithium transport number by a magnesium one and the
magnesium side is the noisier of the two.

The practical reading is narrow and worth stating plainly. A recovery or a permselectivity
quoted from one of these runs should be quoted next to that run's gaps. Five of sixteen
runs balance lithium inside $\pm 10\%$ and three balance magnesium that well. The report
prints both gap columns for every run for this reason, and a number taken out of that
table without them is being quoted more precisely than it was measured.
